% Part: lambda-calculus
% Chapter: introduction
% Section: composition

\documentclass[../../../include/open-logic-section]{subfiles}

\begin{document}

\olfileid{lam}{int}{com}
\olsection{\usetoken{S}{lambda definable} فنکشن ترکیب ہِٹھ بند نیں}

\begin{lem}
!!{lambda definable} فنکشن ترکیب ہِٹھ بند نیں۔
\end{lem}

\begin{proof}
فرض کرو $f$ دی تعریف $h$، $g_0,$
\dots،~$g_{k-1}$ دی ترکیب نال ہوئی اے۔ منّ لئیے کہ $h$، $g_0$،
\dots،~$g_{k-1}$ بالترتیب $H$، $G_0$، \dots،~$G_{k-1}$ نال
!!{lambda defined} نیں۔ سانوں ٹرم~$F$ لبھنا اے جیہڑا~$f$ نوں
!!{lambda define}s۔ اسیں~$F$ نوں سیدھا ایویں متعین کر سکدے آں:
\[
F(x_0, \dots, x_{l-1}) = H(G_0(x_0, \dots, x_{l-1}),
\dots, G_{k-1}(x_0, \dots, x_{l-1})).
\]
دوجے لفظاں وچ لامبڈا کیلکولس دی زبان فنکشوناں دی ترکیب ظاہر کرن
لئی موزوں اے۔
\end{proof}

\end{document}
